% Part: computability
% Chapter: computability-theory
% Section: ce-closed-cup-cap

\documentclass[../../../include/open-logic-section]{subfiles}

\begin{document}

\olfileid{cmp}{thy}{clo}

\olsection[પરિગણનીય ગણોનો યોગગણ અને છેદગણ]{સંગણકીય રીતે પરિગણનીય
  ગણો યોગગણ અને છેદગણ હેઠળ બંધ છે}

આગળનું પ્રમેય સંગણકીય રીતે પરિગણનીય ગણોના કેટલાક બંધત્વના
ગુણધર્મો આપે છે.

\begin{thm}
માનો કે $A$ અને $B$ સંગણકીય રીતે પરિગણનીય છે. તો $A \cap B$
અને $A \cup B$ પણ એવા જ છે.
\end{thm}

\begin{proof}
\olref[eqc]{thm:ce-equiv} વડે સંગણકીય રીતે પરિગણનીય ગણોના
વિવિધ લક્ષણો વાપરી શકીએ છીએ. વાત સ્પષ્ટ કરવા કેટલીક જુદી
સાબિતીઓ આપીશું.

પહેલી સાબિતી માટે માનો કે સંગણનીય વિધેય~$f$, $A$ની અને
સંગણનીય વિધેય~$g$, $B$ની પરિગણના કરે છે. માનો કે
\begin{align*}
h(x) & = \umin{y}{(f(y) = x \lor g(y) = x)} \text{ અને}\\
j(x) & = \umin{y}{(f((y)_0) = x \land g((y)_1) = x)}.
\end{align*}
ત્યારે $A \cup B$, $h$નો પ્રદેશ છે અને $A \cap B$,
$j$નો પ્રદેશ છે.

\begin{explain}
સંગણનાની દૃષ્ટિએ અહીં આવું બને છે: $A$ અને $B$ની પરિગણના
કરતી પ્રક્રિયાઓ હોય, તો કોઈ તત્ત્વ $x$, $A \cup B$માં છે
કે નહીં તેનો અર્ધનિર્ણય બંનેમાંથી એક પરિગણનામાં $x$ને
શોધીને કરી શકાય; અને $x$, $A \cap B$માં છે કે નહીં તેનો
અર્ધનિર્ણય બંને પરિગણનાઓમાં એકસાથે $x$ને શોધીને કરી શકાય.
\end{explain}

બીજી સાબિતી માટે ફરી માનો કે $f$, $A$ની અને $g$, $B$ની
પરિગણના કરે છે. માનો કે
\[
k(x) = \begin{cases}
f(x/2) & \text{જો $x$ સમ હોય} \\
g((x-1)/2) & \text{જો $x$ વિષમ હોય.}
\end{cases}
\]
ત્યારે $k$, $A \cup B$ની પરિગણના કરે છે: $k$ વારાફરતી
$f$ અને~$g$ની પરિગણનાઓમાંથી મૂલ્યો લે છે. $A \cap B$ની
પરિગણના કરવી વધુ મુશ્કેલ છે. $A \cap B$ ખાલી હોય, તો તે
સહેલાઈથી સંગણકીય રીતે પરિગણનીય છે. અન્યથા $A \cap B$નું
કોઈ તત્ત્વ $c$ લો અને $l$ને આ રીતે વ્યાખ્યાયિત કરો:
\[
l(x) = \begin{cases}
f((x)_0) & \text{જો $f((x)_0) = g((x)_1)$} \\
c & \text{અન્યથા.}
\end{cases}
\]
સંગણનાની દૃષ્ટિએ, $l$, $f$ અને $g$ની પરિગણનાઓમાં આવતાં
તત્ત્વોની જોડીઓ તપાસે છે અને મળતો દરેક મેળ આપે છે; અન્યથા
$c$ ફરી આપવાથી પરિગણનામાં કોઈ નવું તત્ત્વ ઉમેરાતું નથી.

છેલ્લી સાબિતી માટે માનો કે $A$, આંશિક વિધેય $m(x)$નો
\emph{પ્રદેશ} છે અને $B$, આંશિક વિધેય $n(x)$નો પ્રદેશ છે.
ત્યારે $A \cap B$, આંશિક વિધેય $m(x) + n(x)$નો પ્રદેશ છે.

\begin{explain}
સંગણનાની દૃષ્ટિએ, $A$ એ એવા મૂલ્યોનો ગણ છે જેના માટે
$m$ થંભે છે અને $B$ એ એવા મૂલ્યોનો ગણ છે જેના માટે
$n$ થંભે છે; $A \cap B$ એવા મૂલ્યોનો ગણ છે જેના માટે
બંને પ્રક્રિયાઓ થંભે છે.
\end{explain}

$A \cup B$ને થંભવાના મૂલ્યોના ગણ તરીકે દર્શાવવું વધુ
મુશ્કેલ છે, કારણ કે $m$ અને $n$ને સમાંતરે અનુકરવા પડે છે.
$d$, $m$નો અને $e$, $n$નો સૂચક હોય; બીજા શબ્દોમાં,
$m = \cfind{d}$ અને $n = \cfind{e}$. ત્યારે $A \cup B$
આ વિધેયનો પ્રદેશ છે:
\[
p(x) = \umin{y}{(T(d,x,y) \lor T(e,x,y))}.
\]
\begin{explain}
સંગણનાની દૃષ્ટિએ, ઇનપુટ $x$ પર $p$, $m$ અથવા $n$ની
થંભતી સંગણના શોધે છે અને બેમાંથી કોઈ એક મળે તો થંભે છે.
\end{explain}
\end{proof}

\end{document}
